\documentclass[aspectratio=169,usenames,dvipsnames]{beamer}

\usetheme{Madrid}
\usecolortheme{default}
\setbeamertemplate{navigation symbols}{}
\setbeamertemplate{footline}[frame number]

\usepackage[utf8]{inputenc}
\usepackage{amsmath}
\usepackage{amssymb}
\usepackage{mathtools}
\usepackage{booktabs}
\usepackage{xcolor}
\usepackage{hyperref}
\usepackage{tikz}
\usepackage{algorithm}
\usepackage{algpseudocode}

\hypersetup{colorlinks=true, linkcolor=blue, urlcolor=blue}

\title{Low-Rank Matrix Factorization}
\subtitle{Comparing Optimization Algorithms for Regularized Matrix Decomposition}
\author{Nicole Chen, Jiuyi (Joy) Cheng, Sirui (Michelle) Li}
\date{Math 347 --- Spring 2026}
\institute{University}

\begin{document}

%=========== SLIDE 1: Title ===========
\begin{frame}
\titlepage
\end{frame}

%=========== SLIDE 2: Project Overview ===========
\begin{frame}{Project Overview \& Roadmap}

\textbf{Project Goal:} Compare three optimization algorithms for the regularized matrix factorization problem:
\[
\min_{U \in \mathbb{R}^{m \times r}, V \in \mathbb{R}^{n \times r}} 
\frac{1}{2}\|UV^T - M\|_F^2 + \frac{\lambda}{2}(\|U\|_F^2 + \|V\|_F^2)
\]

\vspace{0.3cm}

\textbf{Three Algorithms:}
\begin{itemize}
  \item \textbf{Gradient Descent + Armijo:} Simple, stable, pedagogical
  \item \textbf{L-BFGS:} Quasi-Newton, fast convergence
  \item \textbf{Alternating Least Squares (ALS):} Block coordinate descent, ideal for sparse data
\end{itemize}

\vspace{0.3cm}

\textbf{Three Datasets:}
\begin{itemize}
  \item \textbf{Synthetic:} $500 \times 200$, rank-5, $\sigma=0.1$ noise
  \item \textbf{MovieLens 100K:} $943 \times 1682$, 93.7\% sparse, real recommender system
  \item \textbf{Superconductivity:} $21263 \times 81$, dense, real-world feature reduction
\end{itemize}

\vspace{0.3cm}

\textbf{Five Core Experiments:} Method comparison, initialization sensitivity, rank selection, regularization impact, optimization landscape

\end{frame}

%=========== SLIDE 3: Why Matrix Factorization ===========
\begin{frame}{Why Matrix Factorization?}

\textbf{Real-World Problem:}
\begin{itemize}
  \item High-dimensional data with missing or noisy entries
  \item User-item ratings (Netflix problem)
  \item Image denoising and compression
  \item Latent feature discovery
\end{itemize}

\vspace{0.5cm}

\textbf{Key Insight:} Many datasets have underlying low-rank structure
\begin{itemize}
  \item User preferences: $\sim 5$--$10$ latent factors
  \item Images: compressible via wavelets (low-rank in transform domain)
  \item Degrees of freedom: $mr + nr$ factors vs.\ $mn$ entries
  \item For $M = 1000 \times 1700$, $r=20$: compress by $10\times$ vs.\ 2M entries
\end{itemize}

\vspace{0.5cm}

\textbf{Low-rank factorization captures signal, discards noise}

\end{frame}

%=========== SLIDE 4: Problem Statement ===========
\begin{frame}{Problem Statement \& Properties}

\textbf{Regularized Low-Rank Matrix Factorization:}
\begin{align}
\min_{U, V} f(U, V) &= \frac{1}{2}\|UV^T - M\|_F^2 + \frac{\lambda}{2}(\|U\|_F^2 + \|V\|_F^2) \label{eq:objective}
\end{align}

\vspace{0.3cm}

\textbf{Problem Properties:}
\begin{itemize}
  \item \textbf{Non-convex} in $(U, V)$ jointly, but \textbf{biconvex} in each factor separately
  \item \textbf{L-smooth}: $\nabla f$ is Lipschitz continuous (enables convergence)
  \item \textbf{Multiple stationary points}: local minima exist
  \item Smoothness constant: $L \leq \sigma_{\max}(M)^2 + \lambda$ \hyperlink{app:smoothness}{[\checkmark]}
\end{itemize}

\vspace{0.3cm}

\textbf{Regularization trade-off:}
\begin{itemize}
  \item $\lambda \to 0$: Fits noise, high test error
  \item $\lambda$ optimal: Balances fitting signal and avoiding overfitting
  \item $\lambda \to \infty$: Underfitting, poor reconstruction
\end{itemize}

\end{frame}

%=========== SLIDE 5: Applications ===========
\begin{frame}{Applications \& Use Cases}

\begin{columns}[T]
\column{0.5\textwidth}

\textbf{Recommender Systems}
\begin{itemize}
  \item Predict missing user ratings
  \item Netflix Prize benchmark
  \item ~5--50 latent factors sufficient
\end{itemize}

\textbf{Image Processing}
\begin{itemize}
  \item Compression: reduce storage
  \item Denoising: noise is high-rank
  \item Super-resolution: matrix completion
\end{itemize}

\column{0.5\textwidth}

\textbf{Latent Feature Discovery}
\begin{itemize}
  \item Uncover hidden patterns
  \item Dimensionality reduction
  \item Interpretability via factors
\end{itemize}

\textbf{Missing Data Completion}
\begin{itemize}
  \item Inpainting applications
  \item Matrix recovery
  \item Scientific data analysis
\end{itemize}

\end{columns}

\vspace{0.5cm}

\textbf{Key Insight:} Low-rank approximation works when data has intrinsic structure

\end{frame}

%=========== SLIDE 6: Why Low-Rank Works ===========
\begin{frame}{Why Low-Rank Approximation Works}

\textbf{Empirical Observation:} Real data often has intrinsic low-rank structure

\vspace{0.3cm}

\textbf{Degrees of Freedom Analysis:}
\begin{itemize}
  \item Full matrix $M$: $mn$ parameters
  \item Low-rank factors: $mr + nr$ parameters
  \item For $m = 943$, $n = 1682$, $r = 20$:
    \begin{itemize}
      \item Full: $1.59 \times 10^6$ entries
      \item Factors: $3.65 \times 10^4$ entries ($\sim 44\times$ compression)
    \end{itemize}
\end{itemize}

\vspace{0.3cm}

\textbf{Generalization Trade-off:}
\begin{itemize}
  \item Low-rank $\Rightarrow$ fewer parameters $\Rightarrow$ better generalization
  \item Captures signal at rank $r^*$, discards noise
  \item Regularization prevents overfitting to noise
\end{itemize}

\vspace{0.3cm}

\textbf{Signal + Noise Model:}
\[
M = U_{\text{true}} V_{\text{true}}^T + \sigma \cdot \text{Noise}
\]
Optimal $r$ and $\lambda$ recover clean signal from noisy observations

\end{frame}

%=========== SLIDE 7: Datasets ===========
\begin{frame}{Datasets Used in Study}

\begin{center}
\begin{tabular}{|l|c|c|c|l|}
\hline
\textbf{Dataset} & \textbf{Type} & \textbf{Size} & \textbf{Sparsity} & \textbf{Purpose} \\
\hline
Synthetic & Dense & $500 \times 200$ & 0\% & Controlled, ground truth \\
\hline
MovieLens 100K & Sparse & $943 \times 1682$ & 93.7\% & Recommender system \\
\hline
Superconductivity & Dense & $21263 \times 81$ & 0\% & Feature reduction \\
\hline
\end{tabular}
\end{center}

\vspace{0.5cm}

\textbf{Synthetic Data:}
\begin{itemize}
  \item Generation: $M = U_{\text{true}} V_{\text{true}}^T + \sigma \cdot \mathcal{N}(0,1)$
  \item Default: rank 5, $\sigma = 0.1$
  \item Advantage: know true solution $\Rightarrow$ measure recovery error
\end{itemize}

\textbf{MovieLens 100K (Real Sparse):}
\begin{itemize}
  \item 943 users, 1682 movies, $\sim 100$K ratings
  \item Train/test split: evaluate RMSE on held-out ratings
\end{itemize}

\textbf{Superconductivity (Real Dense):}
\begin{itemize}
  \item 21263 compounds, 81 standardized features
  \item Dense matrix: all entries present
\end{itemize}

\end{frame}

%=========== SLIDE 8: Dataset Characteristics ===========
\begin{frame}{Dataset Characteristics Summary}

\begin{center}
\renewcommand{\arraystretch}{1.3}
\begin{tabular}{|c|c|c|c|c|}
\hline
\textbf{Dataset} & \textbf{Type} & \textbf{Size} & \textbf{Sparsity} & \textbf{Est. Rank} \\
\hline
Synthetic & Dense & $500 \times 200$ & 0\% & 5 \\
\hline
MovieLens & Sparse & $943 \times 1682$ & 93.7\% & 20--30 \\
\hline
Superconductivity & Dense & $21263 \times 81$ & 0\% & $\sim 20$ \\
\hline
\end{tabular}
\end{center}

\vspace{0.5cm}

\textbf{Key Observations:}
\begin{itemize}
  \item Synthetic: controlled, ground truth rank known
  \item MovieLens: highly sparse $\Rightarrow$ ALS advantage
  \item Superconductivity: large dense $\Rightarrow$ L-BFGS advantage
\end{itemize}

\end{frame}

%=========== SLIDE 9: Method 1 - Gradient Descent ===========
\begin{frame}{Method 1: Gradient Descent with Armijo Backtracking}

\textbf{Update Rule:}
\[
U_{k+1} = U_k - \alpha_k \nabla_U f(U_k, V_k), \quad 
V_{k+1} = V_k - \alpha_k \nabla_V f(U_k, V_k)
\]

\textbf{Gradients:} \hyperlink{app:grad_dense}{[\checkmark]}
\[
\nabla_U f = (UV^T - M)V + \lambda U, \quad
\nabla_V f = (UV^T - M)^T U + \lambda V
\]

\textbf{Armijo Line Search:}
\begin{itemize}
  \item Start with $\alpha = 1$
  \item If sufficient decrease not met: $\alpha \leftarrow \rho \alpha$ (halve)
  \item Continue until: $f(\theta - \alpha \nabla f) \leq f(\theta) - c \alpha \|\nabla f\|^2$
\end{itemize}

\vspace{0.3cm}

\textbf{Convergence:} \hyperlink{app:gd_convergence}{[\checkmark]} Linear ($\alpha$-linear) to stationary point

\vspace{0.3cm}

\textbf{Advantages \& Disadvantages:}
\begin{itemize}
  \item \textcolor{green}{$\oplus$} Simple, intuitive, stable
  \item \textcolor{red}{$\ominus$} Slow convergence, hyperparameter tuning
\end{itemize}

\end{frame}

%=========== SLIDE 10: Method 2 - L-BFGS ===========
\begin{frame}{Method 2: L-BFGS (Limited-memory BFGS)}

\textbf{Quasi-Newton Method:} Approximate Hessian with limited history
\begin{itemize}
  \item Maintains $m$ recent gradient/step pairs (typically $m = 10$)
  \item Approximates Hessian: $B_k \approx \nabla^2 f$
  \item Update: $x_{k+1} = x_k - \alpha_k B_k^{-1} \nabla f(x_k)$
\end{itemize}

\vspace{0.3cm}

\textbf{Implementation (scipy):}
\begin{itemize}
  \item Pack factors: $x = [U; V] \in \mathbb{R}^{(m+n)r}$
  \item Call \texttt{scipy.optimize.minimize(method='L-BFGS-B')}
  \item Automatic backtracking and Hessian approximation
\end{itemize}

\vspace{0.3cm}

\textbf{Convergence:} \hyperlink{app:lbfgs_convergence}{[\checkmark]} Superlinear: $\|x_k - x^*\| = o(\|x_{k-1} - x^*\|)$
\begin{itemize}
  \item Practical: 2--3$\times$ fewer iterations than GD
\end{itemize}

\vspace{0.3cm}

\textbf{Advantages \& Disadvantages:}
\begin{itemize}
  \item \textcolor{green}{$\oplus$} Fast convergence, few hyperparameters
  \item \textcolor{red}{$\ominus$} More complex, memory overhead
\end{itemize}

\end{frame}

%=========== SLIDE 11: Method 3 - ALS ===========
\begin{frame}{Method 3: Alternating Least Squares (ALS)}

\textbf{Block Coordinate Descent:} Optimize $U$ and $V$ alternately

\vspace{0.3cm}

\textbf{Fix $V$, optimize $U$ (closed form):}
\[
\min_U \frac{1}{2}\|UV^T - M\|_F^2 + \frac{\lambda}{2}\|U\|_F^2
\quad \Rightarrow \quad
U = MV(V^TV + \lambda I)^{-1}
\]

\textbf{Fix $U$, optimize $V$ (symmetric):}
\[
V = M^T U(U^T U + \lambda I)^{-1}
\]

\vspace{0.3cm}

\textbf{Sparse Implementation:} \hyperlink{app:als_sparse}{[\checkmark]}
\begin{itemize}
  \item Row-by-row: for each row $i$, solve ridge regression using only observed entries
  \item Natural handling of missing data
  \item No dense matrix multiplications
\end{itemize}

\vspace{0.3cm}

\textbf{Convergence:} \hyperlink{app:als_convergence}{[\checkmark]} Sublinear descent, guaranteed to stationary point

\vspace{0.3cm}

\textbf{Advantages \& Disadvantages:}
\begin{itemize}
  \item \textcolor{green}{$\oplus$} Closed-form updates, natural for sparse, ideal for real-time
  \item \textcolor{red}{$\ominus$} May converge to different local minima
\end{itemize}

\end{frame}

%=========== SLIDE 12: Method Comparison ===========
\begin{frame}{Method Comparison Summary}

\begin{center}
\renewcommand{\arraystretch}{1.4}
\begin{tabular}{|c|c|c|c|}
\hline
\textbf{Method} & \textbf{Convergence} & \textbf{Per-Iter Cost} & \textbf{Sparse} \\
\hline
GD & Linear ($\alpha$-linear) & $O(mn + r^2)$ & Limited \\
\hline
L-BFGS & Superlinear & $O((m+n)r)$ & Possible \\
\hline
ALS & Sublinear & $O(nnz \cdot r^2)$ & \textbf{Ideal} \\
\hline
\end{tabular}
\end{center}

\vspace{0.5cm}

\textbf{Theoretical Analysis:} \hyperlink{app:convergence}{[\checkmark]}
\begin{itemize}
  \item GD: guaranteed convergence with step size $\alpha \leq 1/L$ (L-smooth constant)
  \item L-BFGS: superlinear convergence from Quasi-Newton theory
  \item ALS: guaranteed descent via biconvexity, sublinear rate typical
\end{itemize}

\vspace{0.3cm}

\textbf{Empirical Trade-offs:}
\begin{itemize}
  \item For \textbf{sparse data}: ALS wins (fewer per-iteration ops)
  \item For \textbf{dense data}: L-BFGS wins (fewer iterations)
  \item For \textbf{stability}: GD most reliable
\end{itemize}

\end{frame}

%=========== SLIDE 13: Experiment 1 - Method Comparison ===========
\begin{frame}{Experiment 1: Method Comparison Across Datasets}

\textbf{Question:} Which algorithm is fastest and most accurate?

\vspace{0.3cm}

\textbf{Setup:}
\begin{itemize}
  \item Run all three methods (GD, L-BFGS, ALS) on all three datasets
  \item Random initialization: $U, V \sim \mathcal{N}(0, 0.1^2)$
  \item Rank: 5 (synthetic), 20 (MovieLens), 20 (Superconductivity)
  \item Regularization: $\lambda = 0.01$
\end{itemize}

\vspace{0.3cm}

\textbf{Metrics:}
\begin{itemize}
  \item Iterations to convergence
  \item Wall-clock time
  \item Final objective value
  \item RMSE (sparse) / Reconstruction error (dense)
\end{itemize}

\vspace{0.3cm}

\textbf{Expected Results:}
\begin{itemize}
  \item ALS: fastest per iteration, especially on MovieLens (sparse)
  \item L-BFGS: fewest iterations overall
  \item GD: slowest but stable
\end{itemize}

\end{frame}

%=========== SLIDE 14: Experiment 2 - Initialization Sensitivity ===========
\begin{frame}{Experiment 2: Initialization Sensitivity}

\textbf{Question:} How sensitive is the problem to initialization?

\vspace{0.3cm}

\textbf{Initializations Tested (21 total configurations):}
\begin{enumerate}
  \item Small random: $U, V \sim \mathcal{N}(0, 0.01^2)$ with 10 seeds
  \item Large random: $U, V \sim \mathcal{N}(0, 1^2)$ with 10 seeds
  \item SVD-based: $U, V$ from truncated SVD of (zero-filled) matrix
\end{enumerate}

\vspace{0.3cm}

\textbf{Key Questions:}
\begin{itemize}
  \item Can poor initialization trap us in bad local minima?
  \item Does SVD initialization provide protection?
  \item Which method most sensitive to initialization?
\end{itemize}

\vspace{0.3cm}

\textbf{Metrics:}
\begin{itemize}
  \item Final objective value
  \item Solution quality (recovery error for synthetic)
  \item Variance across random seeds
\end{itemize}

\vspace{0.3cm}

\textbf{Hypothesis:} SVD initialization should outperform random initialization significantly

\end{frame}

%=========== SLIDE 15: Experiment 3 - Rank Selection ===========
\begin{frame}{Experiment 3: Rank Selection}

\textbf{Question:} How does rank choice affect optimization and reconstruction quality?

\vspace{0.3cm}

\textbf{Rank Grid per Dataset:}
\begin{itemize}
  \item Synthetic: $[2, 3, 5, 8, 10, 20]$ (true rank = 5)
  \item MovieLens: $[2, 5, 10, 20, 30, 50]$
  \item Superconductivity: $[5, 10, 20, 30, 50]$
\end{itemize}

\vspace{0.3cm}

\textbf{Metrics Evaluated:}
\begin{itemize}
  \item Training objective at convergence
  \item Reconstruction error / RMSE
  \item Singular value decay (synthetic only)
\end{itemize}

\vspace{0.3cm}

\textbf{Expected Behavior:}
\begin{itemize}
  \item Underestimating rank: high training error, poor reconstruction (underfitting)
  \item Overestimating rank: potential noise fitting with sparse data (overfitting)
  \item Optimal rank matches intrinsic signal dimension
  \item U-shaped error curve expected
\end{itemize}

\end{frame}

%=========== SLIDE 16: Experiment 4 - Regularization Sensitivity ===========
\begin{frame}{Experiment 4: Regularization Sensitivity}

\textbf{Question:} How does $\lambda$ affect optimization and recovery under noise?

\vspace{0.3cm}

\textbf{Dataset:} Noisy synthetic ($\sigma = 0.5$, higher noise than standard $\sigma = 0.1$)

\vspace{0.3cm}

\textbf{Parameter Grid:}
\[
\lambda \in [10^{-4}, 10^{-3}, 10^{-2}, 10^{-1}, 1.0]
\]

\vspace{0.3cm}

\textbf{Metrics Tracked:}
\begin{itemize}
  \item Final training objective: $f(U, V)$
  \item Recovery error: $\|UV^T - M_{\text{clean}}\|_F / \|M_{\text{clean}}\|_F$
  \item Relative reconstruction error on noisy matrix
\end{itemize}

\vspace{0.3cm}

\textbf{Expected Findings:}
\begin{itemize}
  \item Small $\lambda$: fits noise $\Rightarrow$ high recovery error to clean matrix
  \item Large $\lambda$: underfitting $\Rightarrow$ high training objective
  \item Optimal $\lambda$ trades off bias-variance, recovers signal under noise
\end{itemize}

\end{frame}

%=========== SLIDE 17: Experiment 5 - Optimization Landscape ===========
\begin{frame}{Experiment 5: Optimization Landscape}

\textbf{Question:} How do nonconvex effects appear in practice?

\vspace{0.3cm}

\textbf{Setup (using GD method):}
\begin{enumerate}
  \item Generate 8 solutions from different random seeds (scale = 1.0)
  \item Each runs GD to convergence on synthetic problem
  \item Identify best and worst final objectives
  \item Interpolate: evaluate objective along segment
\end{enumerate}

\vspace{0.3cm}

\textbf{Interpolation Formula:}
\[
f(t) = f((1-t)U_{\text{best}} + tU_{\text{worst}}, (1-t)V_{\text{best}} + tV_{\text{worst}}) 
\quad \text{for } t \in [0, 1]
\]

\vspace{0.3cm}

\textbf{Observable Nonconvex Behavior:}
\begin{itemize}
  \item Non-monotonic objective curve $\Rightarrow$ nonconvexity
  \item Saddle points or local minima along interpolation
  \item Quality variance: $\Delta f = f(U_{\text{worst}}) - f(U_{\text{best}})$
  \item Multiple local minima exist but quality varies
\end{itemize}

\end{frame}

%=========== SLIDE 18: Key Findings - Part 1 ===========
\begin{frame}{Key Findings: Nonconvexity \& Optimization}

\textbf{1. Nonconvexity Doesn't Kill Optimization}
\begin{itemize}
  \item All three methods converge reliably (to possibly different local minima)
  \item With good initialization, convergence is fast
  \item ALS especially robust via block coordinate descent
\end{itemize}

\vspace{0.3cm}

\textbf{2. Rank \& Initialization Matter Most}
\begin{itemize}
  \item Underestimating rank: underfitting bias dominates
  \item Overestimating rank: overfitting risk (especially sparse data)
  \item SVD initialization $\gg$ random initialization
  \item Initialization choice $>$ algorithm choice for solution quality
\end{itemize}

\vspace{0.3cm}

\textbf{3. Regularization Critical for Sparse Data}
\begin{itemize}
  \item MovieLens: $\lambda \in [0.1, 1.0]$ works well
  \item Overly large $\lambda$ kills performance
  \item Dense synthetic: less sensitive to $\lambda$ choice
\end{itemize}

\end{frame}

%=========== SLIDE 19: Key Findings - Part 2 ===========
\begin{frame}{Key Findings: Method Trade-offs}

\textbf{Algorithm Trade-offs:}

\vspace{0.3cm}

\begin{columns}[T]
\column{0.33\textwidth}
\textbf{ALS}
\begin{itemize}
  \item Best for sparse data
  \item Fastest per iteration
  \item Good for real-time systems
  \item Natural missing-data handling
\end{itemize}

\column{0.33\textwidth}
\textbf{L-BFGS}
\begin{itemize}
  \item Fewest iterations
  \item Good for dense offline computation
  \item Robust convergence
  \item Few hyperparameters
\end{itemize}

\column{0.33\textwidth}
\textbf{GD}
\begin{itemize}
  \item Most stable
  \item Easiest to debug
  \item Good pedagogical choice
  \item Predictable behavior
\end{itemize}
\end{columns}

\vspace{0.5cm}

\textbf{Practical Guidance:}
\begin{itemize}
  \item For \textbf{real sparse systems}: ALS with SVD initialization
  \item For \textbf{image/dense tasks}: L-BFGS with SVD initialization
  \item For \textbf{research/teaching}: GD for interpretability
\end{itemize}

\end{frame}

%=========== SLIDE 20: Practical Recommendations ===========
\begin{frame}{Practical Recommendations}

\textbf{For Sparse Data (e.g., Recommender Systems):}
\begin{itemize}
  \item Use ALS with SVD initialization
  \item Regularization: $\lambda = 0.1$ to $1.0$
  \item Rank: 20--50 for typical collaborative filtering
  \item Expect: fast convergence, natural handling of missing data
\end{itemize}

\vspace{0.3cm}

\textbf{For Dense Data (e.g., Image Compression):}
\begin{itemize}
  \item L-BFGS preferred (fewer iterations, fast wall-clock)
  \item SVD initialization essential
  \item Rank determined by compression requirement
  \item Expect: superlinear convergence
\end{itemize}

\vspace{0.3cm}

\textbf{Managing Nonconvexity:}
\begin{itemize}
  \item Initialize well (SVD beats random by large margin)
  \item Regularize appropriately (prevents noise fitting)
  \item Well-regularized problem $\Rightarrow$ few local minima of similar quality
  \item Multiple runs less necessary if initialized well
\end{itemize}

\end{frame}

%=========== SLIDE 21: Conclusions ===========
\begin{frame}{Conclusions}

\textbf{Main Lessons:}

\vspace{0.3cm}

1. \textbf{Low-Rank Factorization is Powerful}
   \begin{itemize}
     \item Applicable across recommender systems, image processing, scientific computing
     \item Compression and denoising benefits clear
   \end{itemize}

\vspace{0.2cm}

2. \textbf{Nonconvex $\neq$ Unsolvable}
   \begin{itemize}
     \item Good initialization + regularization $\Rightarrow$ reliable optimization
     \item Multiple local minima exist but are similar quality
   \end{itemize}

\vspace{0.2cm}

3. \textbf{Algorithm Choice Depends on Data Structure}
   \begin{itemize}
     \item Sparse: ALS best
     \item Dense: L-BFGS best
     \item Stability: GD best
   \end{itemize}

\vspace{0.2cm}

4. \textbf{Open Questions for Future Work}
   \begin{itemize}
     \item Automatic rank selection (cross-validation, information criteria)
     \item Theoretical guarantees avoiding local minima (restrictive conditions needed)
     \item Scalability to billion-scale matrices (stochastic, distributed methods)
   \end{itemize}

\end{frame}

%=========== SLIDE 22: References ===========
\begin{frame}{References \& Resources}

\textbf{Key Papers:}
\begin{itemize}
  \item Zhou et al. (2016): Review of Alternating Least Squares
  \item Candès \& Recht (2009): Matrix completion via convex optimization
  \item Beck \& Teboulle (2009): First-order methods convergence theory
  \item Nesterov (2004): L-smooth functions and gradient descent
\end{itemize}

\vspace{0.3cm}

\textbf{Software Libraries:}
\begin{itemize}
  \item Python: TensorFlow, PyTorch, scikit-learn, scipy.optimize
  \item MATLAB: built-in sparse matrix operations
\end{itemize}

\vspace{0.3cm}

\textbf{Project Implementation:}
\begin{itemize}
  \item GitHub: \texttt{math347}
  \item Code: \texttt{methods.py}, \texttt{experiments.py}, \texttt{objectives.py}
  \item Data: Synthetic, MovieLens 100K, Superconductivity
\end{itemize}

\vspace{0.5cm}

\centering
\Large \textbf{Questions?}

\end{frame}

%=========== APPENDIX ===========
\appendix

\begin{frame}[label=app:smoothness]
\frametitle{Appendix: L-Smoothness}

\textbf{Definition:} Function $f$ is $L$-smooth if:
\[
\|\nabla f(x) - \nabla f(y)\| \leq L \|x - y\| \quad \forall x, y
\]

\vspace{0.3cm}

\textbf{Our Objective:}
\[
f(U, V) = \frac{1}{2}\|UV^T - M\|_F^2 + \frac{\lambda}{2}(\|U\|_F^2 + \|V\|_F^2)
\]

\vspace{0.3cm}

\textbf{Smoothness Constant:}
\[
L \approx \sigma_{\max}(M)^2 + \lambda
\]

\vspace{0.3cm}

\textbf{Consequence:} Gradient descent with step size $\alpha \leq 1/L$ converges to stationary point with rate:
\[
f(x_k) - f^* \leq \mathcal{O}(1/k)
\]

\hyperlink{12}{[\textbf{Back to Slide 12}]}

\end{frame}

\begin{frame}[label=app:grad_dense]
\frametitle{Appendix: Gradient Derivation (Dense Case)}

\textbf{Objective:}
\[
f(U, V) = \frac{1}{2}\|UV^T - M\|_F^2 + \frac{\lambda}{2}(\|U\|_F^2 + \|V\|_F^2)
\]

\vspace{0.3cm}

\textbf{Let} $R = UV^T - M$ (residual matrix):
\[
\|R\|_F^2 = \text{trace}(R^T R) = \text{trace}((UV^T - M)^T(UV^T - M))
\]

\vspace{0.3cm}

\textbf{Gradient w.r.t. $U$:}
\[
\nabla_U f = \frac{\partial}{\partial U}\left[\frac{1}{2}\|R\|_F^2 + \frac{\lambda}{2}\|U\|_F^2\right] = RV + \lambda U = (UV^T - M)V + \lambda U
\]

\vspace{0.3cm}

\textbf{Gradient w.r.t. $V$:}
\[
\nabla_V f = R^T U + \lambda V = (UV^T - M)^T U + \lambda V
\]

\vspace{0.3cm}

\textbf{Computational Cost:} $O(mnr + r^2 m + r^2 n) = O(mnr)$ for full matrix

\hyperlink{9}{[\textbf{Back to Slide 9}]}

\end{frame}

\begin{frame}[label=app:gd_convergence]
\frametitle{Appendix: GD Convergence Analysis}

\textbf{Theorem (Gradient Descent for L-smooth Functions):}

For $L$-smooth convex $f$, GD with step size $\alpha \in (0, 2/L]$ satisfies:
\[
f(x_k) - f(x^*) \leq \frac{\|x_0 - x^*\|^2}{2\alpha k}
\]

\textbf{Rate:} $\mathcal{O}(1/k)$ (linear in $1/k$, i.e., $\alpha$-linear)

\vspace{0.3cm}

\textbf{Our Problem:} Non-convex but $L$-smooth
\begin{itemize}
  \item Guaranteed: $f(x_k) \to$ stationary point
  \item Not guaranteed: $x_k \to$ global optimum
  \item Empirically: often converges to good local minima with good initialization
\end{itemize}

\vspace{0.3cm}

\textbf{Armijo Line Search:} Ensures sufficient decrease, allows adaptive $\alpha$
\[
f(x - \alpha d) \leq f(x) - c \alpha \|\nabla f\|^2, \quad c \in (0, 1)
\]

\hyperlink{9}{[\textbf{Back to Slide 9}]}

\end{frame}

\begin{frame}[label=app:lbfgs_convergence]
\frametitle{Appendix: L-BFGS Convergence}

\textbf{Quasi-Newton Method:} Approximate Hessian $B_k \approx \nabla^2 f$

\vspace{0.3cm}

\textbf{Update:}
\[
x_{k+1} = x_k - B_k^{-1} \nabla f(x_k)
\]

\vspace{0.3cm}

\textbf{L-BFGS (limited-memory):} Store only $m$ recent $(s_i, y_i)$ pairs where:
\begin{itemize}
  \item $s_i = x_{i+1} - x_i$ (step)
  \item $y_i = \nabla f(x_{i+1}) - \nabla f(x_i)$ (gradient change)
\end{itemize}

\vspace{0.3cm}

\textbf{Convergence Theorem (Dennis-Moré):}

L-BFGS with sufficient decrease satisfies:
\[
\|x_k - x^*\| = o(\|x_{k-1} - x^*\|)
\]

\textbf{Rate:} Superlinear convergence

\vspace{0.3cm}

\textbf{Practical:} 2--3$\times$ fewer iterations than GD

\hyperlink{10}{[\textbf{Back to Slide 10}]}

\end{frame}

\begin{frame}[label=app:als_convergence]
\frametitle{Appendix: ALS Convergence}

\textbf{Block Coordinate Descent:} Update $U$ and $V$ in alternation

\vspace{0.3cm}

\textbf{Key Property - Biconvexity:}
\begin{itemize}
  \item $f(U, V)$ non-convex jointly
  \item But: $f$ convex in $U$ (fixing $V$), convex in $V$ (fixing $U$)
  \item Each block update is optimal for fixed other block
\end{itemize}

\vspace{0.3cm}

\textbf{Guaranteed Descent:}
\[
f(U^{(k+1)}, V^{(k)}) \leq f(U^{(k)}, V^{(k)})
\]
\[
f(U^{(k+1)}, V^{(k+1)}) \leq f(U^{(k+1)}, V^{(k)})
\]

Thus $f$ monotonically decreases: $\{f(U^{(k)}, V^{(k)})\}$ is decreasing sequence.

\vspace{0.3cm}

\textbf{Convergence:} To stationary point of $f$ (possibly local minimum)

\textbf{Rate:} Typically sublinear $\mathcal{O}(1/k)$, can accelerate with Nesterov

\hyperlink{11}{[\textbf{Back to Slide 11}]}

\end{frame}

\begin{frame}[label=app:als_sparse]
\frametitle{Appendix: Sparse ALS Implementation}

\textbf{Sparse Update for Row $i$ of $U$:}

Fix $V$, solve for $u_i$:
\[
\min_{u_i} \frac{1}{2}\sum_{j \in \Omega_i} (u_i^T v_j - M_{ij})^2 + \frac{\lambda}{2}\|u_i\|^2
\]

where $\Omega_i = \{j : M_{ij} \text{ is observed}\}$

\vspace{0.3cm}

\textbf{Solution (Ridge Regression):}
\[
u_i = \left(\sum_{j \in \Omega_i} v_j v_j^T + \lambda I\right)^{-1} \sum_{j \in \Omega_i} M_{ij} v_j
\]

\vspace{0.3cm}

\textbf{Sparse Update for Column $j$ of $V$:}
\[
v_j = \left(\sum_{i \in \Omega_j} u_i u_i^T + \lambda I\right)^{-1} \sum_{i \in \Omega_j} M_{ij} u_i
\]

where $\Omega_j = \{i : M_{ij} \text{ is observed}\}$

\vspace{0.3cm}

\textbf{Computational Benefit:} $O(|\Omega_i| r^2)$ per row instead of $O(nr^2)$ for dense case

\hyperlink{11}{[\textbf{Back to Slide 11}]}

\end{frame}

\begin{frame}[label=app:convergence]
\frametitle{Appendix: Convergence Rates Summary}

\begin{center}
\renewcommand{\arraystretch}{1.8}
\begin{tabular}{|l|c|c|}
\hline
\textbf{Method} & \textbf{Convergence Rate} & \textbf{Analysis} \\
\hline
\textbf{GD} & $\mathcal{O}(1/k)$ ($\alpha$-linear) & L-smooth function \\
\hline
\textbf{L-BFGS} & Superlinear & Quasi-Newton, Dennis-Moré \\
\hline
\textbf{ALS} & $\mathcal{O}(1/k)$ (sublinear) & Biconvex block descent \\
\hline
\end{tabular}
\end{center}

\vspace{0.5cm}

\textbf{Detailed References:}

\begin{itemize}
  \item \textbf{GD:} Nesterov (2004), "Introductory Lectures on Convex Optimization"
  \item \textbf{L-BFGS:} Nocedal (1980), "Updating quasi-Newton matrices with limited storage"
  \item \textbf{ALS:} Zhou et al. (2016), "Alternating Least Squares for Low-Rank Matrix Reconstruction"
\end{itemize}

\vspace{0.3cm}

\textbf{Note:} Rates are worst-case asymptotic. Practical convergence often faster with good initialization.

\hyperlink{12}{[\textbf{Back to Slide 12}]}

\end{frame}

\end{document}
